% Generated draft: requires mathematical and source review.
\documentclass[11pt]{article}
\usepackage[margin=25mm]{geometry}
\usepackage{fontspec}
\setmainfont{Latin Modern Roman}
\usepackage{amsmath,amssymb,mathtools,unicode-math}
\setmathfont{Latin Modern Math}
\usepackage{enumitem,needspace,xurl,hyperref}
\setlength{\parindent}{0pt}
\setlength{\parskip}{5pt}
\setlength{\emergencystretch}{3em}
\newcommand{\N}{\mathbb{N}}
\newcommand{\Z}{\mathbb{Z}}
\newcommand{\Q}{\mathbb{Q}}
\newcommand{\R}{\mathbb{R}}
\newcommand{\C}{\mathbb{C}}
\newcommand{\F}{\mathbb{F}}
\newcommand{\A}{\mathbb{A}}
\newcommand{\PP}{\mathbb{P}}
\newcommand{\E}{\mathbb{E}}
\newcommand{\eps}{\varepsilon}
\newcommand{\dd}{\,\mathrm{d}}
\begin{document}
\section*{UnsolvedMath ID: 30002049}
\textbf{Source problem code:} catalogue entry
\subsection{Definitions and mathematical statement}
The supplied record does not specify the ground field or its characteristic. To express its assertion, let \(k\) denote an unresolved algebraically closed ground field and let \(N\geq 1\) be the ambient dimension. This notation does not assert that the conjecture is intended over every such field.

Write \(\mathbb G_{m,k}^{N}\) for the algebraic torus with \(k\)-points \((k^{\times})^{N}\). A subtorus is a connected algebraic subgroup isomorphic to \(\mathbb G_{m,k}^{d}\) for some \(0\leq d\leq N\). A torsion point \(\zeta\) is a point satisfying \(\zeta^{m}=(1,\ldots,1)\) for some integer \(m\geq 1\), with powers taken coordinatewise. A torsion coset is a translate \(\zeta H\), where \(H\) is a subtorus and \(\zeta\) is a torsion point. All codimensions below are taken in \(\mathbb G_{m,k}^{N}\), so that \(\operatorname{codim}T=N-\dim T\). Dimension means algebraic dimension; for a nonempty reducible algebraic subset it is the maximum dimension of its irreducible components.

\textbf{Supplied conjectural assertion (attributed to Zilber).} For every algebraic subset \(W\subseteq\mathbb G_{m,k}^{N}\), there exists a finite collection \(\mathcal U_W\) of torsion cosets of codimension one in \(\mathbb G_{m,k}^{N}\) such that, for every torsion coset \(T_0\subseteq\mathbb G_{m,k}^{N}\) and every irreducible component \(Y\) of \(W\cap T_0\),
\[
\dim Y>\dim W-\operatorname{codim}T_0
\quad\Longrightarrow\quad
\exists\,T\in\mathcal U_W\text{ such that }Y\subseteq W\cap T.
\]
The family \(\mathcal U_W\) is chosen for \(W\) before quantifying over \(T_0\) and \(Y\). The statement imposes no irreducibility or equidimensionality assumption on \(W\); its inequality uses \(\dim W\), exactly as supplied. When \(W\) is empty, the condition is vacuous and the empty family suffices. No effectiveness requirement is included in this assertion. The intended ground field and characteristic require clarification before this becomes a fully specified mathematical statement.
\subsection{Short English statement}
For each algebraic subset of a torus, the conjecture asks for finitely many torsion translates of codimension-one subtori that contain every intersection component whose dimension exceeds the supplied expected-dimension bound. The intended ground field is unspecified. The supplied record labels the conjecture open; this draft does not independently verify that status.
\subsection{Sources}
\begin{itemize}
\item Local UnsolvedMath registry, permanent ID \(30002049\), title \emph{Zilber Conjecture for Torsion-Anomalous Intersections}, exact source problem code \texttt{catalogue entry}, in the Oberwolfach Reports Open Problems collection. The record attributes the conjecture to Zilber, lists proposed year 2012, and names no individual proposer.
\item Supplied report citation: \emph{Diophantische Approximationen (2013)}, Oberwolfach Reports, DOI 10.4171/owr/2012/22, pp. 1307--1359. \url{https://doi.org/10.4171/owr/2012/22}. The background describes these reports as documenting mathematical workshops and their open-problem discussions. The supplied statement-verification annotation says that the official report gives this formulation with the ambient torus and exceptional family explicit; that annotation has not been independently verified here.
\item Amoroso--Leroux--Sombra, \emph{Overdetermined systems of sparse polynomial equations}. \url{https://arxiv.org/abs/1307.5788}.
\item Capuano, \emph{Exploring Zilber's conjecture on atypical intersections}. \url{https://doi.org/10.1112/jlms.70536}.
\item Supplied literature assessment dated exactly \texttt{2026-08-21}: the record labels the problem open and describes it as Zilber's unlikely-intersection conjecture for algebraic tori. It states that the cited sparse-equation algorithm is conditional on an effective form, and that the general conjecture remains open despite many special cases. This is supplied dated curation derived from the OWR material, not an independently verified literature assessment. The record expressly calls for independent verification of time-sensitive status and cited claims.
\end{itemize}
\end{document}
